\section{5. Conclusion}


In this paper, we propose a Unified Dual-Mask Physical Model for non-Lambertian light field depth estimation to overcome the severe geometric corruption caused by complex reflectance properties. By bridging physical optics with epipolar plane image (EPI) (EPI) analysis, our framework fundamentally departs from the conventional Lambertian photo-consistency assumption, enabling robust material disentanglement and accurate depth recovery in real-world mixed-material scenes. 

First, we establish a dual-mask physical model via MRI-inspired angular frequency analysis, which unifies material perception across Lambertian, non-Lambertian, and mixed scenes to yield reliable physical priors and high-precision material distribution maps. Second, we devise a component-aware depth estimation strategy based on inverse wavevector analysis, fundamentally resolving the inherent failure of traditional EPI-based methods in specular regions and preventing prediction degradation into mere texture replication. Third, we formulate a physics-prior-driven three-stage progressive training framework that guarantees physical interpretability, circumvents overfitting during black-box optimization, and substantially enhances generalization capabilities on small-sample non-Lambertian datasets. 

Although our approach achieves superior depth accuracy and robustness, it currently relies on densely sampled angular views, which restricts its direct deployment on sparse light field arrays or low-cost plenoptic cameras with limited angular resolution. Future implementation will focus on extending the inverse wavevector analysis to sparse angular sampling paradigms by integrating adaptive view-synthesis priors, thereby ensuring reliable geometric reasoning under severe angular sparsity.


The proposed angular frequency parsing framework fundamentally reframes light field depth estimation from a spatial-domain matching problem to an angular-frequency filtering task. By treating the $9 \times 9$ angular sampling matrix as a discrete spatial signal, the 2D Discrete Fourier Transform (DFT) exploits the convolution theorem to decouple the incident illumination from the surface reflectance. In the angular frequency domain, the medium mask and angular direction mask act as adaptive bandpass filters. Specifically, the diffuse component concentrates its energy at the DC coefficient (zero frequency), while the specular and scattering components manifest as distinct high- and mid-frequency spectral lobes, respectively. This spectral separation mathematically validates the dual-mask assumption, demonstrating that material disentanglement can be achieved through linear phase shifts in the $k$-space without requiring non-linear, data-driven feature extraction.

From an engineering deployment perspective, the transition to the angular frequency domain yields substantial computational dividends. Conventional deep learning approaches for non-Lambertian depth estimation typically rely on 4D convolutions or volumetric cost volumes, incurring $\mathcal{O}(U^2 V^2 HW)$ complexity and prohibitive memory footprints. In contrast, our component-aware strategy leverages the Fast Fourier Transform (FFT) and pixel-wise mask operations, strictly bounding the computational complexity to $\mathcal{O}(HW \log(UV))$, which simplifies to $\mathcal{O}(HW)$ for a fixed angular resolution. This algorithmic lightweight nature ensures that the physical priors can be integrated into real-time pipelines for autonomous navigation and robotic manipulation, where strict latency constraints and limited onboard compute preclude the use of heavy black-box neural networks.

Despite its robustness, the physical model encounters theoretical boundaries under extreme optical conditions. When surface roughness approaches zero (i.e., perfect mirrors), the angular spectrum degenerates into a discrete Dirac delta function. In such cases, the high-frequency wave vectors become highly susceptible to angular aliasing if the sampling rate violates the Nyquist-Shannon criterion, potentially leading to erroneous mask activation. Furthermore, the current dual-mask formulation assumes local light-surface interactions governed by standard microfacet theories. It does not inherently account for non-local light transport phenomena, such as deep subsurface scattering in translucent materials (e.g., wax, marble, or human skin) or complex multi-bounce inter-reflections in highly concave geometries. In these scenarios, the angular frequency distribution becomes heavily smeared and spatially shifted, no longer strictly adhering to the predefined mid-frequency scattering lobe, which may cause the inverse wavevector analysis to underestimate depth discontinuities.

Finally, the frequency-domain analysis provides actionable insights for the design of next-generation light field hardware. The explicit mapping between material roughness and angular spectral bandwidth implies that a uniform angular sampling grid is inherently suboptimal. To faithfully capture the high-frequency wave vectors of highly specular regions without aliasing, the micro-lens array must satisfy a localized, material-dependent angular Nyquist rate. This suggests a paradigm shift toward adaptive or non-uniform angular sampling in future plenoptic cameras. By dynamically allocating angular resolution based on the high-frequency energy distribution in the optical $k$-space, hardware designers can maximize the information entropy of the captured light field under a fixed sensor pixel budget, directly bridging the gap between physical optical modeling and computational sensor design.

Despite the demonstrated efficacy of the unified dual-mask physical model in disentangling complex reflectance properties, several limitations remain from a signal processing and physical optics perspective, outlining promising directions for future work.

First, the current angular frequency-domain parsing relies on a 2D discrete Fourier transform (2D-DFT) over a localized $9 \times 9$ angular grid. While this efficiently captures the primary energy distributions for isotropic diffuse, specular, and scattering components, it inherently assumes local spatial stationarity. For highly anisotropic materials, such as brushed metals or woven fabrics, the bidirectional reflectance distribution function (BRDF) exhibits direction-dependent high-frequency variations. In the angular $k$-space, these manifest as asymmetric, non-linear energy manifolds that cannot be fully resolved by standard 2D-DFT magnitude spectra. Future work will investigate higher-order spectral analysis, such as fractional Fourier transforms or localized wavelet packets, to better characterize anisotropic wavevector fields. Additionally, incorporating spherical harmonic expansions into the angular frequency analysis could provide a more rigorous mathematical representation of direction-dependent scattering lobes, thereby refining the directional mask estimation.

Second, the medium mask currently quantifies pixel-level surface roughness to dictate light-matter interaction types. However, this formulation struggles with multi-layered dielectric materials exhibiting subsurface scattering or clear-coat effects, where incident light undergoes multiple internal reflections before exiting. In such scenarios, the superposition of distinct wavevectors corrupts the pure surface reflectance priors. To address this, we plan to integrate polarized light field imaging into our physical model. Polarization cues provide orthogonal physical constraints that can effectively decouple surface specularities from subsurface diffuse scattering, extending the dual-mask framework to handle complex volumetric compositions. Furthermore, coupling the medium mask with micro-facet distribution models could allow for the explicit recovery of physically meaningful roughness parameters rather than relative qualitative metrics.

Finally, our inverse wavevector analysis is strictly formulated for static scenes. In dynamic environments, the temporal evolution of non-Lambertian reflections introduces severe aliasing in the spatiotemporal frequency domain, breaking the static wavevector continuity assumptions. Extending the proposed MRI-inspired $k$-space analogy to a 3D spatiotemporal frequency volume will be crucial. By designing temporal anti-aliasing filters and dynamic wavevector tracking mechanisms, the framework could achieve robust depth estimation for non-Lambertian dynamic scenes. This extension would preserve the $O(HW)$ spatial complexity while ensuring temporal coherence, making the system viable for real-time robotic manipulation in unstructured environments.